\section{Put the Page with the Values}
\label{sec:ch13-put-page-with-values}
\index{entity reference!from search results}%

Search returns references to Sessions rather than copying canonical titles into
its public contract. Its local Session shape is only an opaque node id:

\begin{minted}{csharp}
using HotChocolate.Types.Relay;

namespace Search;

public sealed record Session([property: ID] string Id);
\end{minted}

The key is not resolvable here. Search can identify a Session, but the router
must ask Sessions to hydrate it.

\begin{minted}{csharp}
using HotChocolate.ApolloFederation.Types;

namespace Search;

[ObjectType<Session>]
[Key("id", false)]
public static partial class SessionType;
\end{minted}

The document resolver formats the integer database key as the same global node
id used everywhere else in the graph:

\begin{minted}{csharp}
using HotChocolate.Types.Relay;

namespace Search;

[ObjectType<SessionSearchDocument>]
public static partial class SessionSearchDocumentType
{
    public static Session GetSession(
        [Parent] SessionSearchDocument document,
        INodeIdSerializer serializer)
        => new(serializer.Format(nameof(Session), document.SessionId));
}
\end{minted}

That produces a composed shape in which discovery data and canonical entity
data remain distinguishable:

\begin{graphqlsdl}
type SessionSearchDocument {
  session: Session!
  startsAt: DateTime!
  averageScore: Float
  ratingCount: Int!
}
\end{graphqlsdl}

I keep the statement interceptor in this service too. Without a separate log,
the claim that paging happened locally would be guesswork.

\begin{minted}{csharp}
using System.Data.Common;
using Microsoft.EntityFrameworkCore.Diagnostics;

namespace Search;

public sealed class StatementLog : DbCommandInterceptor
{
    private static int _count;

    public override InterceptionResult<DbDataReader> ReaderExecuting(
        DbCommand command,
        CommandEventData eventData,
        InterceptionResult<DbDataReader> result)
    {
        Console.WriteLine(
            $"-- statement {Interlocked.Increment(ref _count)}");
        Console.WriteLine(command.CommandText);
        return result;
    }

    public override
        ValueTask<InterceptionResult<DbDataReader>> ReaderExecutingAsync(
        DbCommand command,
        CommandEventData eventData,
        InterceptionResult<DbDataReader> result,
        CancellationToken cancellationToken = default)
    {
        Console.WriteLine(
            $"-- statement {Interlocked.Increment(ref _count)}");
        Console.WriteLine(command.CommandText);
        return ValueTask.FromResult(result);
    }
}
\end{minted}

The host is the same explicit stack used by the other subgraphs. Paging
arguments are enabled before federation, cost defaults remain disabled for the
composer, and the database is created from the deterministic seed.

\begin{minted}{csharp}
using ChilliCream.Nitro.App;
using HotChocolate.AspNetCore;
using Microsoft.EntityFrameworkCore;
using Search;

var builder = WebApplication.CreateBuilder(args);

builder.Logging.AddFilter(
    "Microsoft.EntityFrameworkCore.Database.Command",
    LogLevel.Warning);

builder.Services.AddDbContextFactory<SearchContext>(options =>
{
    options.UseSqlite("Data Source=search.db");
    if (builder.Environment.IsDevelopment())
    {
        options.AddInterceptors(new StatementLog());
    }
});

builder
    .AddGraphQL()
    .AddSearchTypes()
    .RegisterDbContextFactory<SearchContext>()
    .AddPagingArguments()
    .AddDefaultNodeIdSerializer()
    .AddApolloFederation()
    .ModifyCostOptions(options => options.ApplyCostDefaults = false);

var app = builder.Build();

using (var scope = app.Services.CreateScope())
{
    var factory = scope.ServiceProvider
        .GetRequiredService<IDbContextFactory<SearchContext>>();
    using var db = factory.CreateDbContext();
    db.Database.EnsureCreated();
}

app.MapGraphQL().WithOptions(options =>
{
    options.Tool.ServeMode = ServeMode.Embedded;
});

app.RunWithGraphQLCommands(args);
\end{minted}

Finally, composition gets the fourth input. This is the complete
\code{graph/graph.yaml}, not a patch:

\begin{minted}[highlightlines={15,16,17,18}]{yaml}
version: 1
subgraphs:
  - name: sessions
    routing_url: http://localhost:5001/graphql
    schema:
      file: ../src/Sessions/schema.graphql
  - name: speakers
    routing_url: http://localhost:5002/graphql
    schema:
      file: ../src/Speakers/schema.graphql
  - name: ratings
    routing_url: http://localhost:5003/graphql
    schema:
      file: ../src/Ratings/schema.graphql
  - name: search
    routing_url: http://localhost:5004/graphql
    schema:
      file: ../src/Search/schema.graphql
\end{minted}

The router config does not change, but a service addition is exactly when I
want to reprint it. There is still one static execution artifact and one
router process:

\begin{minted}{yaml}
version: "1"

# Without this the router asks Cosmo Cloud for its execution config and
# refuses to start, because it has no token to ask with. The path is
# resolved against the directory the router is started in, not against
# this file, so start it from this folder:
#
#   router -config config.yaml
execution_config:
  file:
    path: ../graph/router.json

# A development switch, and what makes the query plan visible at all.
# The plan is an Advanced Request Tracing option, and outside dev mode
# the router wants a Cosmo Cloud token before it hands one out. With
# this off, X-WG-Include-Query-Plan is ignored rather than refused.
dev_mode: true

# Who the router thinks is asking. The alternative entry in this list
# is a jwks url, which is what a deployment with an identity provider
# behind it uses; the shared secret is here because this book stands
# up no such provider. header_key_id is not optional: the router
# matches it against the kid in the token header, so the token this
# book mints carries one.
#
# A request carrying a token this key does not verify is answered 401
# by the router and never reaches a subgraph.
authentication:
  jwt:
    jwks:
      - secret: splitting-the-graph-development-signing-key-32b
        symmetric_algorithm: HS256
        header_key_id: dev

# The router does not forward request headers to subgraphs unless it
# is told to, and a subgraph that guards its own data cannot tell who
# is asking without this block. Leaving it out does not make the graph
# insecure; it makes the guarded field unreachable. Every request that
# names one comes back null with an authorization error underneath,
# for a caller holding a good token exactly as for a caller holding
# none, which is what makes the missing block read like a policy
# rather than like a header. A request that never names a guarded
# field is unaffected, so this survives a smoke test.
headers:
  all:
    request:
      - op: propagate
        named: Authorization
\end{minted}

The fourth subgraph changes graph ownership, not the router's configuration
model.
